\section*{Chapitre \ref{ch:g11:fossil-fuel-geology} --- La géologie du charbon, du pétrole et du gaz}

\begin{solution}{exo:g11:fossil-fuel-geology:1}
Le plancton mort se dépose sur un fond marin pauvre en oxygène~; la
matière organique devient du \omterm{def:g11:fossil-fuel-geology:source-rock}{kérogène} dans une boue~; l'enfouissement sous
des sédiments plus jeunes~; la \omterm{def:g11:fossil-fuel-geology:oil-window}{maturation} dans la \omterm{def:g11:fossil-fuel-geology:oil-window}{fenêtre à huile}~; la
\omterm{def:g11:fossil-fuel-geology:migration}{migration}~; le pétrole s'accumule sous une \omterm{def:g11:fossil-fuel-geology:reservoir}{couverture}.
\end{solution}

\begin{solution}{exo:g11:fossil-fuel-geology:2}
(1) \omterm{def:g11:fossil-fuel-geology:source-rock}{Roche mère}. (2) \omterm{def:g11:fossil-fuel-geology:reservoir}{Roche réservoir}. (3) \omterm{def:g11:fossil-fuel-geology:reservoir}{Couverture}. (4) \omterm{def:g11:fossil-fuel-geology:reservoir}{Couverture}.
(5) \omterm{def:g11:fossil-fuel-geology:reservoir}{Roche réservoir}.
\end{solution}

\begin{solution}{exo:g11:fossil-fuel-geology:3}
Sommet~: $(60 - 10)/25 = \qty{2.0}{km}$~; base~:
$(120 - 10)/25 = \qty{4.4}{km}$. La fenêtre est épaisse de
$60/25 = \qty{2.4}{km}$.
\end{solution}

\begin{solution}{exo:g11:fossil-fuel-geology:4}
(1) Tourbe, lignite, \omterm{def:g11:fossil-fuel-geology:coal}{charbon} bitumineux, anthracite. (2) L'anthracite, de
86 à \qty{97}{\percent} de carbone. (3) L'anthracite, au-delà d'environ
\qty{5}{km} et \qty{150}{\degreeCelsius}.
\end{solution}

\begin{solution}{exo:g11:fossil-fuel-geology:5}
(1) Les fluides se rangent par densité~: le gaz est le plus léger, puis le
pétrole, puis l'eau. (2) Le pétrole et le gaz montent parce qu'ils sont plus
légers que l'eau~; seule une roche qu'ils ne peuvent pas traverser les
arrête. La pélite (ou l'argilite) et le sel.
\end{solution}

\begin{solution}{exo:g11:fossil-fuel-geology:6}
(1) A~: $(60 - 15)/20 = \qty{2.25}{km}$ et $(120 - 15)/20 = \qty{5.25}{km}$.
B~: $45/40 \approx \qty{1.13}{km}$ et $105/40 \approx \qty{2.63}{km}$.
(2) À \qty{2}{km}~: dans A, $15 + 40 = \qty{55}{\degreeCelsius}$, la roche
n'est pas encore mature et ne produit pas de pétrole~; dans B,
$15 + 80 = \qty{95}{\degreeCelsius}$, dans la \omterm{def:g11:fossil-fuel-geology:oil-window}{fenêtre à huile}, elle produit
du pétrole. (3) Dans le bassin A~: le pétrole ne s'y forme qu'au-dessous de
\qty{2.25}{km}.
\end{solution}

\begin{solution}{exo:g11:fossil-fuel-geology:7}
(1) $11 \times 3 = \qty{33}{m}$ de tourbe. (2) $\qty{33000}{mm}$ à raison
de \qty{1}{mm} par an~: environ \num{33000} ans. (3) La tourbe ne se
conserve que sous une eau stagnante, hors d'atteinte de l'oxygène~; à
mesure qu'elle s'accumule, le sol doit s'enfoncer au même rythme pour que
sa surface reste gorgée d'eau.
\end{solution}

\begin{solution}{exo:g11:fossil-fuel-geology:8}
(1) $20 + 30 \times 1.2 = \qty{56}{\degreeCelsius}$~;
$20 + 30 \times 2.3 = \qty{89}{\degreeCelsius}$~;
$20 + 30 \times 4.0 = \qty{140}{\degreeCelsius}$.
(2) Il y a 120 millions d'années, immature~; il y a 95 millions d'années,
dans la \omterm{def:g11:fossil-fuel-geology:oil-window}{fenêtre à huile}~; il y a 20 millions d'années, dans la \omterm{def:g11:fossil-fuel-geology:oil-window}{fenêtre à
gaz}.
\end{solution}

\begin{solution}{exo:g11:fossil-fuel-geology:9}
(1) $2580 - 2400 = \qty{180}{m}$. (2) Le gaz occupe
$2450 - 2400 = \qty{50}{m}$, le pétrole $2580 - 2450 = \qty{130}{m}$.
(3) Ils s'échappent au point de fuite, sous la courbe fermée la plus
profonde, et poursuivent leur \omterm{def:g11:fossil-fuel-geology:migration}{migration} vers le haut, jusqu'à un autre
\omterm{def:g11:fossil-fuel-geology:trap}{piège} ou jusqu'à la surface.
\end{solution}

\begin{solution}{exo:g11:fossil-fuel-geology:10}
(1) Sommet de la fenêtre à \qty{3}{km}~: $(60 - 20)/3 \approx
\qty{13}{\degreeCelsius/km}$~; base~: $(120 - 20)/3 \approx
\qty{33}{\degreeCelsius/km}$. (2) À \qty{3}{km}, la température va de
$20 + 75 = \qty{95}{\degreeCelsius}$ à $20 + 102 = \qty{122}{\degreeCelsius}$~:
dans la \omterm{def:g11:fossil-fuel-geology:oil-window}{fenêtre à huile}, ou à sa base pour les gradients les plus forts.
La règle est raisonnable~: dans la plupart des bassins, une \omterm{def:g11:fossil-fuel-geology:source-rock}{roche mère}
enfouie sous \qty{3}{km} de roches a commencé à produire du pétrole.
\end{solution}

\begin{solution}{exo:g11:fossil-fuel-geology:11}
(1) \qty{30}{\percent}~: lignite~; \qty{70}{\percent}~: \omterm{def:g11:fossil-fuel-geology:coal}{charbon}
bitumineux~; \qty{92}{\percent}~: anthracite. (2) L'anthracite, au-delà
d'environ \qty{5}{km} et \qty{150}{\degreeCelsius}. (3) Noir et vitreux.
\end{solution}

\begin{solution}{exo:g11:fossil-fuel-geology:12}
(1) $1100/25 = \qty{44}{m}$ par million d'années. (2) $0.044 \times 30
\approx \qty{1.3}{\degreeCelsius}$ par million d'années. (3) La fenêtre
s'étend sur \qty{60}{\degreeCelsius}~: $60/1.32 \approx 45$ millions
d'années. (4) L'enfouissement a ralenti après 95 millions d'années (environ
24, puis 23 et \qty{20}{m} par million d'années)~: la \omterm{def:g11:fossil-fuel-geology:source-rock}{roche mère} s'est
réchauffée plus lentement.
\end{solution}

\begin{solution}{exo:g11:fossil-fuel-geology:13}
(1) $(80 - 51)/(117 - 51) = 29/66 \approx 0.44$~: environ
\qty{44}{\percent}. (2) Il a migré vers le haut avant qu'aucun \omterm{def:g11:fossil-fuel-geology:trap}{piège}
n'existe au-dessus~: il s'est échappé vers la surface, où il a été oxydé, ou
s'est accumulé dans un autre \omterm{def:g11:fossil-fuel-geology:trap}{piège}. (3) Il y a 40 millions d'années, la
\omterm{def:g11:fossil-fuel-geology:source-rock}{roche mère} avait quitté la \omterm{def:g11:fossil-fuel-geology:oil-window}{fenêtre à huile}~: elle ne produisait plus que du
gaz, et l'anticlinal n'aurait pu en retenir que du gaz.
\end{solution}

\begin{solution}{exo:g11:fossil-fuel-geology:14}
(1) Sommet~: $(60 - 20)/40 = \qty{1.0}{km}$~; base~: $(120 - 20)/40 =
\qty{2.5}{km}$. (2) \qty{1.0}{km} se situe entre \qty{0.4}{km} (il y a 150
millions d'années) et \qty{1.2}{km} (120)~: $150 - 30 \times 0.6/0.8 = 127.5$ millions d'années. \qty{2.5}{km} se situe entre \qty{2.3}{km} (95)
et \qty{3.0}{km} (66)~: $95 - 29 \times 0.2/0.7 \approx 86.7$ millions
d'années. (3) Pas de pétrole~: la \omterm{def:g11:fossil-fuel-geology:source-rock}{roche mère} a quitté la \omterm{def:g11:fossil-fuel-geology:oil-window}{fenêtre à huile}
avant la formation de l'anticlinal, et son pétrole avait migré ailleurs.
Il pourrait retenir le gaz produit après 80 millions d'années.
\end{solution}

\begin{solution}{exo:g11:fossil-fuel-geology:15}
(1) Au moins \qty{5}{km}. (2) $15 + 30 \times 5 = \qty{165}{\degreeCelsius}$,
au-dessus de \qty{150}{\degreeCelsius}~: c'est cohérent. (3) Au moins
\qty{5}{km} de roches ont été soulevés et érodés. $\qty{5e6}{mm}$ en au
plus \num{299} millions d'années donnent au moins $5 \times 10^6/2.99 \times 10^8
\approx \qty{0.017}{mm}$
par an en moyenne.
\end{solution}

\begin{solution}{pb:g11:fossil-fuel-geology:1}
\textbf{1.} La \omterm{def:g11:fossil-fuel-geology:source-rock}{roche mère}~: l'argilite noire~; le réservoir~: le grès
poreux~; la \omterm{def:g11:fossil-fuel-geology:reservoir}{couverture}~: la pélite qui le surmonte.

\textbf{2.} Le plancton déposé sur un fond marin pauvre en oxygène a
échappé à la décomposition et a été enfoui dans la boue.

\textbf{3.} La \omterm{def:g11:fossil-fuel-geology:source-rock}{roche mère} au Jurassique (165 se situe entre 201.4 et
143.1). Le grès a commencé à se déposer au Jurassique, mais
essentiellement au Crétacé inférieur.

\textbf{4.} La matière organique dispersée dans une roche sédimentaire,
insoluble dans les solvants organiques. Elle ne produit du pétrole qu'une
fois chauffée vers \qty{60}{\degreeCelsius}, ce qui demande un
enfouissement.

\textbf{5.} $400 + 800 + 1100 + 700 + 700 + 300 - 500 = \qty{3500}{m}$.

\textbf{6.} \qty{0.4}{km}, \qty{1.2}{km}, \qty{2.3}{km}, \qty{3.0}{km},
\qty{3.7}{km}, \qty{4.0}{km}~; aujourd'hui \qty{3.5}{km}.

\textbf{7.} $400/15 \approx 26.7$~; $800/30 \approx 26.7$~; $1100/25 = 44$~;
$700/29 \approx 24.1$~; $700/31 \approx 22.6$~; $300/15 = 20$ (\unit{m} par
million d'années).

\textbf{8.} Il y a 20 millions d'années, à \qty{4.0}{km}.

\textbf{9.} Le bassin a été soulevé et \qty{500}{m} de roches ont été
érodés à son sommet.

\textbf{10.} La base du grès s'est déposée il y a 150 millions d'années,
sur \qty{400}{m} d'argilite et de pélite~: $3.5 - 0.4 = \qty{3.1}{km}$.

\textbf{11.} $T = 20 + 30 z$~: \qty{32}{\degreeCelsius}, \num{56}, \num{89},
\num{110}, \num{131}, \qty{140}{\degreeCelsius}~; aujourd'hui
\qty{125}{\degreeCelsius}.

\textbf{12.} $40/30 \approx \qty{1.33}{km}$ et $100/30 \approx
\qty{3.33}{km}$.

\textbf{13.} Entre il y a 120 millions d'années (\qty{56}{\degreeCelsius})
et il y a 95 millions d'années (\qty{89}{\degreeCelsius}).

\textbf{14.} La température a augmenté de 33 degrés en 25 millions
d'années~; il lui a fallu $25 \times 4/33 \approx 3.0$ millions d'années
pour gagner les 4 degrés de \num{56} à \qty{60}{\degreeCelsius}~: il y a
environ $120 - 3.0 = 117$ millions d'années.

\textbf{15.} Entre il y a 66 millions d'années (\qty{110}{\degreeCelsius})
et il y a 35 (\qty{131}{\degreeCelsius})~: il y a
$66 - 31 \times 10/21 \approx 51.2$ millions d'années.

\textbf{16.} \qty{140}{\degreeCelsius}, il y a 20 millions d'années. À
\qty{125}{\degreeCelsius} aujourd'hui, elle est dans la \omterm{def:g11:fossil-fuel-geology:oil-window}{fenêtre à gaz}~:
elle produit surtout du méthane.

\textbf{17.} Sans \omterm{def:g11:fossil-fuel-geology:trap}{piège} au-dessus, il a migré vers le haut et s'est
échappé, vers la surface ou vers d'autres \omterm{def:g11:fossil-fuel-geology:trap}{pièges}.

\textbf{18.} $(80 - 51.2)/(117.0 - 51.2) \approx 0.44$.

\textbf{19.} Du gaz en haut, puis du pétrole, puis de l'eau.

\textbf{20.} La \omterm{def:g11:fossil-fuel-geology:source-rock}{roche mère} est entrée dans la \omterm{def:g11:fossil-fuel-geology:oil-window}{fenêtre à huile} il y a
environ \textbf{117 millions d'années}.
\end{solution}
